% Response to Reviewer Comments — Applied Sciences
% Manuscript: Parametric Investigation of Thermal Characteristics
%             in High-Speed Angular Contact Ball Bearings
%             via Transient Multi-Physics Simulation
%
% Original title: ADORE: A High-Fidelity Six-Degree-of-Freedom Transient
%                 Thermo-Elasto-Hydrodynamic Bearing Dynamics Solver
%                 with Closed-Loop Energy Conservation
%
% Note: The manuscript has been substantially revised and refocused as a
% parametric thermal analysis paper, as described below.

\documentclass[11pt]{article}
\usepackage[margin=2.5cm]{geometry}
\usepackage{xcolor}
\usepackage{enumitem}
\usepackage{hyperref}

\definecolor{response}{RGB}{0,80,160}

\newcommand{\reviewer}[1]{\subsection*{#1}}
\newcommand{\comment}[1]{\textbf{Comment:} \textit{#1}\par\medskip}
\newcommand{\response}[1]{\textbf{Response:} \textcolor{response}{#1}\par\bigskip}

\title{Point-by-Point Response to Reviewer Comments}
\author{}
\date{}

\begin{document}
\maketitle

\noindent We thank both reviewers for their thorough and constructive evaluations. The feedback has substantially improved the manuscript. In this revision, we have \textbf{refocused the paper} from a solver-description paper to a \textbf{parametric thermal analysis} paper, which better matches the scope of Applied Sciences and directly addresses the reviewers' concerns about over-claiming solver generality. The title has been changed accordingly.

Below we address each comment point by point. Blue text indicates our response; cross-references to sections and tables refer to the revised manuscript.

\bigskip
\noindent\rule{\textwidth}{0.4pt}

%% =====================================================================
\section*{Reviewer 1}
%% =====================================================================

\reviewer{Major Comment 1: Validation claims stronger than evidence (120-mm case)}

\comment{The cross-bearing prediction for the 120-mm case shows a systematic underprediction of about 70\%. The abstract and conclusions should distinguish clearly between the calibrated 35-mm case and the unsuccessful blind extrapolation to 120~mm.}

\response{We fully agree. In the revised manuscript, the 120-mm cross-prediction has been moved from the Validation section to the Discussion (\S5.4, ``Cross-Bearing Size Limitation''), where it is explicitly framed as a \emph{model limitation} rather than a validation success. The abstract and conclusions now scope all claims to the 35-mm bearing only. The ``high-fidelity'' and ``first-principles'' terminology has been removed throughout.}


\reviewer{Major Comment 2: Calibration and validation not separated}

\comment{The 35-mm NASA case appears to be used for parameter calibration. The manuscript should separate calibration from validation and state the digitization uncertainty.}

\response{The traction curve parameters (Table~3) are taken from the bearing dynamics literature [Gupta 1984, Harris 2006] and are \emph{not} fitted to the NASA 35-mm thermal data. This is now stated explicitly in \S2.3 (line~118 of the revised manuscript): ``the parameter values are representative of polyol-ester lubricants and are taken from the ranges reported in the bearing dynamics literature---they are \emph{not} fitted to the NASA 35-mm thermal data used for validation.''

Additionally, a digitization uncertainty of $\pm 3\%$ for the NASA reference data has been added to \S3.1 (line~345).}


\reviewer{Major Comment 3: Four-node vs.\ five-node thermal inconsistency}

\comment{The manuscript alternates between four-node and five-node descriptions. This must be corrected throughout.}

\response{The revised manuscript consistently uses ``four-node'' throughout (abstract, \S2.4, \S3, conclusions, captions). The ambient temperature is a boundary condition, not a dynamic state. A thorough consistency audit of the thermal network size, state vector definition, and state-count formula has been performed.}


\reviewer{Major Comment 4: Clearance/preload treatment insufficiently defined}

\comment{Initial clearance/preload state and its role in contact geometry are not defined clearly. A sensitivity analysis showing how predictions respond to clearance/preload variation is needed.}

\response{This is now fully addressed:
\begin{itemize}[nosep]
    \item Equations~5--8 (\S2.4.1) define thermal expansion of all three components (inner race, outer race, ball).
    \item Equation~8 defines the effective clearance including centrifugal effects.
    \item A clearance sensitivity sweep (0--30~$\mu$m) at the nominal 667~N axial preload is presented in \S4.3 (Figures~6--7).
    \item A supplementary sweep at low preload ($F_a = 100$~N, $P_d$ = 0--100~$\mu$m) was performed. Both sweeps show zero thermal sensitivity, which is physically correct under full preload and is acknowledged as a model limitation for clearance-dependent load-zone effects (\S4.3, Limitation paragraph).
\end{itemize}}


\reviewer{Major Comment 5: Reproducibility claim not fully supported}

\comment{Several quantities are not fully linked to the tabulated parameter set. A complete mapping from symbols to values is needed.}

\response{The revised manuscript includes six complete parameter tables (Tables~1--5) covering geometry, lubricant, traction curve, thermal network (with both physical and accelerated values), and drag/churning parameters. All symbolic variables in the equations are now traceable to specific table entries. The FBDF integrator is named, tolerances (\texttt{rtol}, \texttt{atol}) and maximum step size are fully specified (\S2.6), and a footnote in Table~3 documents the source of the traction parameters.}


\reviewer{Major Comment 6: Thermal acceleration weakens temperature claims}

\comment{The thermal capacitances were reduced by 100$\times$. This changes the thermal time scale and cannot be treated as direct validation.}

\response{To address this concern, we performed a full physical-time simulation (1$\times$ capacitance, 3.0~s duration) alongside the accelerated case (100$\times$, 0.15~s) at the baseline operating condition. The results (Table~6, Figure~4) show that all steady-state temperatures agree within 0.2~K and total heat generation within 0.09\%. This confirms that the thermal acceleration technique does not alter the equilibrium point. We have also added a rigorous caveat that this equivalence is guaranteed only when the system possesses a unique stable fixed point (\S4.1.1, line~457).}


\reviewer{Major Comment 7: Additional independent validation needed}

\comment{The paper would be strengthened by comparison with experimental temperature data, friction torque measurements, or dynamic signatures.}

\response{The revised manuscript adds:
\begin{itemize}[nosep]
    \item Quantitative friction torque comparison with Palmgren and SKF models at four speed points (Table~7, \S3.2), showing the simulation consistently falls between the two bounds.
    \item Predicted steady-state temperatures at four speeds (Table~8, \S3.3) with ranking consistent with Takabi and Khonsari~(2013) experimental observations.
    \item A new literature comparison subsection (\S5.3) discussing consistency with published temperature differentials.
\end{itemize}
We acknowledge that direct per-node temperature validation remains unavailable for this bearing and note this as a limitation (\S3.5).}


\reviewer{Minor Comments}

\response{All minor comments have been addressed:
\begin{itemize}[nosep]
    \item MDPI Applied Sciences template is now used throughout.
    \item Author affiliation is finalized.
    \item Terminology standardized (four-node throughout).
    \item Over-categorical claims softened; all conclusions scoped to ``this specific bearing configuration.''
    \item NASA data digitization uncertainty ($\pm 3\%$) stated in \S3.1.
    \item Figures reformatted with informative captions including units and operating conditions.
\end{itemize}}


\bigskip
\noindent\rule{\textwidth}{0.4pt}

%% =====================================================================
\section*{Reviewer 2}
%% =====================================================================

\reviewer{Major Comment 1: Four/five-node thermal inconsistency}

\response{Addressed identically to Reviewer~1 Major Comment~3. The manuscript now consistently describes a four-node LPTN throughout all sections.}


\reviewer{Major Comment 2: 120-mm cross-bearing reveals major limitation}

\comment{The 120-mm prediction shows 70\% underprediction. The section should be reframed as a model limitation with quantitative sensitivity tests.}

\response{The cross-bearing analysis has been moved to the Discussion (\S5.4) and is now explicitly labeled as a ``model limitation.'' We identify three contributing factors: increased churning losses, size-dependent traction scaling, and different thermal boundary conditions. While a full quantitative decomposition is reserved for future work, the consistent $3\times$ offset provides a clear indicator of the model's validity envelope.}


\reviewer{Major Comment 3: ``First-principles'' claim should be moderated}

\comment{The traction parameters appear to be selected differently for different cases. The ``first-principles'' label should be moderated.}

\response{The terms ``first-principles'' and ``high-fidelity'' have been removed entirely from the revised manuscript. The framework is now described as a ``validated transient simulation tool.'' The traction parameters (Table~3) are explicitly stated to be from the published literature [Gupta 1984, Harris 2006], not fitted to the NASA 35-mm thermal data (\S2.3).}


\reviewer{Major Comment 4: Validation limited in independence}

\comment{Harris cage speed is not experimental validation; energy conservation is internal consistency. At least one additional independent validation source is needed.}

\response{The revised paper refocuses validation on three independent lines of evidence: (1)~NASA TP-2275 heat generation (Table~5), (2)~Palmgren/SKF friction torque bounding (Table~7), and (3)~Takabi \& Khonsari temperature trend consistency (\S5.3). The Harris cage speed comparison and energy audit are no longer presented as primary validation; energy conservation is an internal solver metric not included in the paper.}


\reviewer{Major Comment 5: Lack of sensitivity analysis}

\comment{The results likely depend strongly on thermal conductances, traction coefficients, oil fill fraction, etc.}

\response{A comprehensive thermal conductance sensitivity analysis has been added (\S3.6, Table~9). All five $G_{ij}$ values are simultaneously perturbed by $\pm 25\%$ and $\pm 50\%$. The results show:
\begin{itemize}[nosep]
    \item Absolute $T_{\text{IR}}$ uncertainty: $-16$ to $+29$~K for $\pm 50\%$ $G$ variation.
    \item Inner--outer race differential remains stable ($5.6$--$6.1$~K).
    \item Relative parametric trends are robust (all sweep cases use the same $G_{ij}$).
\end{itemize}
Additionally, a LPTN circumferential limitation estimate (\S3.6.1) quantifies the peak ball temperature error at 3--5~K under asymmetric loading.}


\reviewer{Major Comment 6: Journal formatting}

\response{The manuscript has been completely reformatted using the MDPI Applied Sciences template (\texttt{mdpi.cls}).}


\reviewer{Minor Comments 1--10}

\response{All minor comments addressed:
\begin{itemize}[nosep]
    \item[\textbf{1.}] Thermal-network terminology: consistently ``four-node'' throughout.
    \item[\textbf{2.}] State-vector: total dimension $N_{\text{state}} = 13Z + 28 = 236$ clearly stated (\S2.1).
    \item[\textbf{3.}] Traction--shear link: the constitutive chain (traction curve $\to$ Bair--Winer $\to$ Archard flash temperature) is described in \S2.3 with explicit citations.
    \item[\textbf{4.}] Figure captions: all revised with units, operating conditions, and physical interpretation.
    \item[\textbf{5.}] Gupta comparison: removed from the revised, refocused manuscript.
    \item[\textbf{6.}] Baumgarte stabilization: not relevant to the revised scope (quaternion kinematics not discussed).
    \item[\textbf{7.}] References: Takabi \& Khonsari now substantively discussed in \S5.3; additional recent references added (Zheng 2018, Hamrock 1977, Bair 1979, Goksem 1978, Rackauckas 2017).
    \item[\textbf{8.}] Churning citation: Goksem and Hargreaves (1978) reference added (\S2.5, Ref.~[16]).
    \item[\textbf{9.}] Code availability: version identifier and README reproduction note added (\S Data Availability).
    \item[\textbf{10.}] FBDF rationale: motivation over IRK/Rosenbrock methods now stated (\S2.6).
\end{itemize}}

\bigskip
\noindent\rule{\textwidth}{0.4pt}
\section*{Summary of Changes}

The manuscript has been substantially revised and refocused:
\begin{enumerate}[nosep]
    \item \textbf{Scope}: Refocused from a solver-description paper to a parametric thermal analysis paper, with the title changed accordingly.
    \item \textbf{Validation}: Three independent validation lines (NASA heat generation, Palmgren/SKF bounding, Takabi temperature trends); 120-mm reframed as limitation.
    \item \textbf{Sensitivity}: $G_{ij}$ $\pm 50\%$ perturbation study (Table~9) and LPTN limitation estimate added.
    \item \textbf{Thermal acceleration}: Validated with full physical-time simulation ($\Delta T < 0.2$~K).
    \item \textbf{Reproducibility}: Six parameter tables, FBDF with tolerances, traction source documented.
    \item \textbf{Clearance}: Full treatment with Eqs.~5--8, two sweep ranges, model limitation acknowledged.
    \item \textbf{Language}: All over-claims removed; conclusions scoped to the 35-mm bearing.
\end{enumerate}

\end{document}
